\section{Consideraciones éticas}
\label{sec:etica}

El estudio no involucra participantes humanos ni datos personales, por lo que
los resguardos éticos no se orientan al consentimiento informado sino a la
responsabilidad profesional en el manejo de información industrial y en la
comunicación de resultados de ingeniería. Los resguardos se estructuran bajo el
esquema riesgo~$\rightarrow$~acción de mitigación.

\subsection{Riesgo 1: Sobreatribución de madurez tecnológica}

\textbf{Riesgo.} Presentar un modelo matemático validado exclusivamente sobre
datos sintéticos como si fuera un producto informático listo para sustentar
decisiones operativas reales en planta. Este es el riesgo central y más propio
de este estudio. Su materialización podría inducir falsas expectativas de
seguridad o de rendimiento industrial, y en un contexto donde las decisiones de
alimentación de un chancador tienen consecuencias sobre la integridad del equipo
y sobre la seguridad de las personas, esa sobreatribución no es un problema
meramente académico.

\textbf{Acción de mitigación.} Se declara de manera explícita, en el informe, en
el artículo científico y en la presentación oral, que el algoritmo propuesto se
encuentra en un nivel de madurez tecnológica de laboratorio y simulación. Se
exige de forma expresa el pilotaje físico \emph{in situ} y la calibración
empírica de seguridad como condiciones previas a cualquier despliegue comercial
en un gemelo digital operativo. Esta declaración se incorpora además al archivo
\texttt{README} del código entregado, de modo que acompañe al artefacto y no
solo al documento.

\subsection{Riesgo 2: Exposición de información industrial confidencial}

\textbf{Riesgo.} Vulnerar los acuerdos de no divulgación vigentes en la
industria minera mediante la exposición de telemetría propietaria, leyes de
mineral o rendimientos comerciales de faenas específicas.

\textbf{Acción de mitigación.} El proyecto prescinde por completo de datos
telemétricos crudos de yacimientos específicos y fundamenta su análisis en datos
sintéticos calibrados con literatura científica pública. Los parámetros
dimensionales utilizados corresponden a rangos genéricos de la gran minería a
cielo abierto documentados en fuentes abiertas, y no a valores de operación de
ninguna faena identificable. Con ello se anula el riesgo de filtración de
información estratégica.

\subsection{Riesgo 3: Integridad de los resultados reportados}

\textbf{Riesgo.} Reportar métricas de desempeño no verificables, ajustadas para
favorecer la conclusión, o no reproducibles por terceros.

\textbf{Acción de mitigación.} Se adopta un principio de reproducibilidad
estricta: todo valor numérico reportado en este informe proviene de la ejecución
del código anexo con semilla fija, y cualquier evaluador puede reejecutarlo y
obtener exactamente los mismos resultados. Se reportan tanto las métricas
favorables como las desfavorables —en particular, el rechazo del supuesto de
normalidad por la prueba de Shapiro-Wilk y las tres combinaciones de
sintonización en que el módulo resulta peor que no filtrar—, y se documentan de
forma explícita las recalibraciones de parámetros realizadas durante el
desarrollo, con su justificación física (sección~\ref{sec:recalibracion}).

\subsection{Riesgo 4: Trazabilidad y veracidad de las fuentes}

\textbf{Riesgo.} Sustentar afirmaciones del marco teórico o del estado del arte
en referencias no verificables, lo que comprometería la integridad académica del
trabajo y la validez de la brecha de conocimiento declarada.

\textbf{Acción de mitigación.} La totalidad de las referencias citadas fue
verificada individualmente contra su fuente primaria antes de la consolidación
de este informe. Se retiraron aquellas cuya existencia no pudo comprobarse y se
sustituyeron por fuentes verificadas de pertinencia equivalente o superior. Se
corrigió además la atribución editorial de una referencia que aparecía asignada
a una revista incorrecta en entregas previas. El Anexo~\ref{anexo:matriz}
consigna, para cada fuente, su identificador digital de objeto, lo que permite
la verificación independiente por parte del evaluador.

\subsection{Riesgo 5: Uso responsable de herramientas de asistencia}

\textbf{Riesgo.} Incorporar contenido generado por herramientas automatizadas
sin verificación, en particular datos, citas o resultados numéricos.

\textbf{Acción de mitigación.} Todo resultado numérico presentado procede de
código ejecutado y auditable, no de estimaciones. Toda referencia bibliográfica
fue contrastada con su fuente. El código, los registros de trazabilidad y las
figuras se entregan como anexo digital, de modo que la cadena completa entre el
dato y la afirmación sea inspeccionable.
